\chapter{Introducción}
\label{cap:introduccion}

\section{Problema}
\label{sec:problema}

\subsection{Contexto operativo}

El monitoreo sismo-volcánico en Chile es responsabilidad del Observatorio
Volcanológico de los Andes del Sur (OVDAS), dependiente de SERNAGEOMIN. Su
labor depende de la capacidad de detectar, clasificar y localizar eventos
sísmicos de forma continua sobre el flujo permanente de datos que producen las
redes desplegadas en cada complejo volcánico. El procesamiento automático no
reemplaza al analista: le permite concentrar su atención en la actividad
relevante en lugar de revisar manualmente el registro completo.

El Complejo Volcánico Nevados de Chillán (NVChVC) presenta cinco tipos de
señal de interés para esa vigilancia. Tres son de origen volcánico:
volcano-tectónicos (VT), de largo período (LP) y tremor (TR). Dos no lo son,
pero aparecen en el mismo registro y con energía comparable: avalanchas (AV) y
eventos de hielo o \emph{ice-quakes} (IC). El sistema debe distinguir unas de
otras, porque solo las primeras informan sobre el estado interno del volcán.

El sistema tomado como línea base en esta tesis es \texttt{ovdas-core}
\parencite{aracena2025ovdascore}, una reimplementación del flujo de trabajo
del observatorio bajo arquitectura de microservicios. Su etapa de detección
emplea el algoritmo STA/LTA \parencite{allen1978stalta} con umbral fijo, y su
etapa de clasificación una red convolucional VGG16 \parencite{simonyan2015vgg}
de tres salidas, acompañada de un módulo de detección de señales fuera de
distribución (\emph{out-of-distribution}, OOD) basado en vecinos más cercanos
que fuerza la etiqueta \enquote{OT} cuando la señal no se parece a ninguna
clase conocida.

\subsection{Limitaciones del enfoque vigente}
\label{subsec:limitaciones}

Esa arquitectura arrastra cuatro limitaciones que no son defectos de
implementación, sino consecuencias de decisiones de diseño que nunca fueron
contrastadas empíricamente sobre datos del volcán objetivo.

La primera es que \textbf{el umbral de detección decide sobre energía, no
sobre clase}. Una avalancha es un transitorio energético de banda ancha y un
\emph{ice-quake} es impulsivo y limpio; ambos disparan un detector de energía
tan bien como un evento volcánico. No existe un valor de umbral que admita los
eventos volcánicos débiles y rechace los no volcánicos intensos, porque las
dos poblaciones se solapan en la variable sobre la que el umbral decide.

La segunda es que \textbf{la elección del detector no está respaldada por
evidencia local}. STA/LTA fue calibrado manualmente y los detectores basados
en aprendizaje profundo disponibles hoy fueron entrenados con catálogos
tectónicos. \textcite{jiang2021comparacion} muestran caídas de \emph{recall}
de entre 13\,\% y 56\,\% al aplicar PhaseNet y EQTransformer fuera de su
región de entrenamiento, de modo que el desempeño publicado no puede asumirse
representativo del contexto sismo-volcánico chileno.

La tercera es que \textbf{la ventana de análisis es de longitud fija}. El
clasificador opera sobre segmentos de 8\,192 muestras, equivalentes a
81,92\,s a 100\,Hz. Esa ventana no proviene de la duración de los eventos sino
de una conveniencia arquitectónica: la red espera una entrada de tamaño
constante.

La cuarta es que \textbf{el clasificador se entrena con recortes de referencia
y opera sobre recortes automáticos}. Los ejemplos de entrenamiento provienen
de segmentaciones hechas por analistas, mientras que en operación los límites
del evento los fija el detector. Nadie ha medido cuánto desempeño se pierde en
ese traspaso, pese a que es la condición real de uso.

\subsection{Evidencia empírica preliminar}
\label{subsec:evidencia}

La fase experimental ya ejecutada en esta tesis convierte esas cuatro
limitaciones en mediciones. Los resultados se obtuvieron sobre un tramo
continuo de 10,18\,h del NVChVC con 203 eventos catalogados utilizables,
proveniente del conjunto de datos público asociado a
\textcite{espinosacurilem2026segmentacion}
\parencite{espinosacurilem2025dataset}, y se detallan en el
Capítulo~\ref{cap:resultados}.

Sobre el solapamiento entre clases de interés y clases descartables, el
catálogo de referencia muestra que \textbf{73 de los 203 eventos, un 36\,\%,
son avalanchas o \emph{ice-quakes}}. No son un caso marginal: son más de un
tercio de lo que el sistema encuentra. Y los detectores no fallan al
encontrarlos, sino al contrario: EQTransformer los detecta significativamente
mejor que PhaseNet, con diferencias de $-0{,}432$ en AV y $-0{,}301$ en IC
cuyos intervalos de confianza excluyen el cero.

Sobre la comparación entre detectores, los tres candidatos resultan
\textbf{estadísticamente indistinguibles al medirlos globalmente}, con
\emph{recall} de 0,542, 0,527 y 0,581 e intervalos de confianza del 95\,\% que
se solapan por completo. El desglose por clase, en cambio, revela diferencias
significativas de signo opuesto que se cancelan al promediar: PhaseNet aventaja
en LP y TR, y EQTransformer en VT, AV e IC, con significancia en las cinco
clases. Reportar un único valor de \emph{recall} global para un detector
sismo-volcánico oculta información esencial sobre su comportamiento.

Sobre la ventana fija, la \textbf{duración mediana de los tremores es de
87,8\,s y excede los 81,92\,s de la ventana}, con un rango que va de 3,4\,s a
240,8\,s. La representación estándar trunca sistemáticamente una clase
completa, y ninguna elección de tamaño fijo puede cubrir un rango de dos
órdenes de magnitud sin sacrificar resolución en un extremo o contexto en el
otro.

Sobre la calidad del recorte, STA/LTA supera a EQTransformer de forma robusta
frente al umbral de emparejamiento (IoU medio 0,563 contra 0,479), mientras
que PhaseNet queda excluido de esa comparación por construcción, ya que
entrega arribos de fase y no intervalos delimitados. El análisis de los 106
falsos positivos de STA/LTA muestra además que un 20,8\,\% son fragmentos de
eventos catalogados y un 35,8\,\% son detecciones adyacentes a ellos: la mayor
parte de lo que se contabiliza como error de detección es, en realidad, un
problema de delimitación.

\subsection{Percepción de las partes interesadas}
\label{subsec:interesados}

La relevancia del problema está documentada por quienes operan y desarrollan
estos sistemas. \textcite{espinosacurilem2025ood}, trabajando sobre el mismo
complejo volcánico, plantean explícitamente que los clasificadores de conjunto
cerrado asignan con alta confianza una clase conocida a señales que no
pertenecen a ninguna de las clases de entrenamiento, y que esa conducta es
inaceptable en un contexto de vigilancia donde una clasificación errónea puede
alterar la lectura del estado del volcán.
\textcite{espinosacurilem2026segmentacion} sitúan el problema en términos
operativos: el objetivo es el reconocimiento en tiempo real sobre registro
continuo, no la clasificación de eventos previamente aislados por un analista.

A esa evidencia documental se suma una medición propia sobre el sistema en
operación: al ejecutar \texttt{ovdas-core} sobre sus archivos de prueba, el
módulo OOD rechaza entre el 59\,\% y el 100\,\% de los eventos detectados según
el archivo, asignándoles la etiqueta \enquote{OT}. Un sistema que descarta esa
proporción de lo que él mismo detecta traslada al analista el trabajo que
debía ahorrarle.

% --------------------------------------------------------
AQUI VA LO DE LAS percepción de las partes interesadas que no tengo idea que es
% ---------------------------------------------------------------

\subsection{Síntesis del problema}
\label{subsec:sintesis-problema}

El problema puede enunciarse así: \textbf{no existe evidencia empírica local
que permita decidir qué detector conviene usar sobre registro sismo-volcánico
del NVChVC, ni un procedimiento de recorte que se ajuste a la duración real de
los eventos, ni una medición de cuánto desempeño pierde un clasificador al
pasar de recortes de referencia a recortes automáticos.} Mientras esas tres
cosas no existan, cualquier mejora del clasificador se construye sobre
supuestos no verificados.

\section{Estado del Arte}
\label{sec:estado-arte}

Esta sección sintetiza los trabajos que sustentan directamente las decisiones
metodológicas de esta tesis, sin pretender una revisión exhaustiva del campo.

\subsection{De los detectores heurísticos al aprendizaje profundo}

El algoritmo STA/LTA \parencite{allen1978stalta} compara la energía media en
una ventana corta con la de una ventana larga y declara una detección cuando
el cociente supera un umbral. Es computacionalmente económico, no requiere
entrenamiento y sigue siendo el detector desplegado en numerosos
observatorios, pero su desempeño depende de la calibración manual de umbrales
y ventanas, y no se adapta a cambios en el régimen de ruido.

\textcite{zhu2019phasenet} presentan PhaseNet, un método de \emph{picking} de
fases P y S basado en una arquitectura U-Net unidimensional que produce, para
cada instante, una distribución de probabilidad sobre arribo P, arribo S y
ruido. Fue entrenado con más de 700\,000 formas de onda del norte de
California, un catálogo puramente tectónico, y reporta F1$_P$ = 0,896 y
F1$_S$ = 0,801 en esa misma región. Es el más liviano de los detectores
profundos considerados, factor relevante bajo la restricción de inferencia en
CPU descrita en la Sección~\ref{sec:alcances}.

\textcite{mousavi2020eqtransformer} presentan EQTransformer, que resuelve
simultáneamente detección y \emph{picking} mediante atención jerárquica sobre
el catálogo STEAD. A diferencia de PhaseNet, entrega un canal de detección que
delimita el evento además de los arribos, lo que resulta determinante para
esta tesis: solo los detectores que delimitan pueden alimentar directamente un
recorte. Ambos modelos se emplean aquí a través de \texttt{seisbench}
\parencite{woollam2022seisbench}, que unifica su interfaz de uso.

\subsection{Generalización fuera del dominio de entrenamiento}

\textcite{jiang2021comparacion} comparan ambos modelos sobre las réplicas de
los terremotos de Yangbi y Maduo, fuera de su región de entrenamiento, y
documentan caídas de \emph{recall} de entre 13\,\% y 56\,\% respecto de lo
publicado. Encuentran además que PhaseNet supera a EQTransformer en ese
contexto pese a su menor complejidad. Este resultado es la justificación
metodológica central del primer objetivo específico de esta tesis: el
desempeño publicado no se asume, se verifica sobre datos locales.
\textcite{sheehan2025eqtlab} refuerzan el punto desde otro ángulo, mostrando
que EQTransformer preentrenado transfiere a emisiones acústicas de
laboratorio, un dominio aún más alejado del tectónico.

\subsection{Clasificación de eventos sismo-volcánicos}

\textcite{bueno2020bnn} aplican redes bayesianas con transferencia de
aprendizaje a seis clases sismo-volcánicas, incorporando cuantificación de
incertidumbre, y alcanzan 92,08\,\% de exactitud.
\textcite{espinosacurilem2025ood} trabajan sobre el mismo complejo volcánico
de esta tesis y combinan un VGG16 con un detector OOD basado en vecinos más
cercanos sobre las representaciones internas de la red; su mejor configuración
combina espectrograma y forma de onda, con F1 media de 95,3\,\% y AUROC de
0,97 para la detección de señales desconocidas. Esta es la arquitectura ya
implementada en \texttt{ovdas-core} y se adopta como línea base de comparación
en esta tesis, no como diseño propio: la limitación que se busca superar es
precisamente que sus clases no volcánicas se resuelven por rechazo estadístico
en lugar de por supervisión, pese a existir ejemplos etiquetados de ellas.

\subsection{Detección y clasificación conjuntas}

\textcite{espinosacurilem2026segmentacion} proponen resolver detección y
clasificación en una sola pasada: apilan las formas de onda de hasta ocho
estaciones en una imagen bidimensional y aplican modelos de segmentación
semántica que asignan una clase por píxel entre fondo y cinco tipos de evento.
Reportan F1 = 0,91 e IoU = 0,88 con UNet sobre ventanas de 80\,s. El enfoque
se evalúa aquí como punto de comparación de resultados, pero no se adopta:
al fundir detección y clasificación en un solo modelo eliminaría la comparación
independiente de detectores, que es el objeto del primer objetivo específico.
El conjunto de datos publicado junto a ese trabajo
\parencite{espinosacurilem2025dataset} es el que se emplea en esta tesis.

\subsection{Posicionamiento de esta tesis}

De la revisión se desprenden tres observaciones que fundamentan el diseño
metodológico. Primero, la generalización geográfica no puede asumirse, por lo
que los tres detectores se ejecutan sobre datos locales antes de decidir.
Segundo, ningún trabajo revisado mide la degradación que sufre un clasificador
al pasar de recortes de referencia a recortes producidos automáticamente, pese
a que esa es la condición de operación real. Tercero, la ventana de longitud
fija es una convención heredada de la visión computacional y no una propiedad
de los datos sísmicos, cuyas duraciones abarcan dos órdenes de magnitud.

La contribución de esta tesis no consiste en proponer arquitecturas de red
nuevas, sino en configurar, adaptar y evaluar componentes existentes para
construir un sistema integrado ajustado a este problema, y en medir lo que
hasta ahora no se ha medido.

\section{Objetivos}
\label{sec:objetivos}

\subsection{Objetivo general}

\begin{itemize}
    \item Desarrollar un algoritmo de identificación de eventos
    sismo-volcánicos mediante detección automática y clasificación con redes
    neuronales convolucionales, ajustado a la variabilidad del recorte
    entregado por el detector, sobre registros sísmicos continuos del Complejo
    Volcánico Nevados de Chillán.
\end{itemize}

\subsection{Objetivos específicos}
\label{subsec:objetivos-especificos}

\begin{itemize}
    \item \textbf{OE1.} Evaluar el desempeño de detectores automáticos de
    eventos sismo-volcánicos en dos niveles, identificación del instante del
    evento y delimitación temporal de su recorte, sobre registros continuos
    etiquetados del Complejo Volcánico Nevados de Chillán.

    \item \textbf{OE2.} Construir un procedimiento de recorte adaptativo y de
    representación espectral de los eventos detectados, ajustado a su duración
    y a su solapamiento, compatible con la entrada de un clasificador
    convolucional.

    \item \textbf{OE3.} Validar el algoritmo integrado de identificación sobre
    el registro continuo, contrastando arquitecturas convolucionales
    alternativas y cuantificando la degradación entre el recorte de referencia
    y el recorte automático.
\end{itemize}

El objetivo general emplea un solo verbo en infinitivo perteneciente al nivel
\emph{crear}, el más alto de la taxonomía de Bloom revisada
\parencite{anderson2001taxonomy}. Los tres objetivos específicos se ubican en
los niveles \emph{evaluar}, \emph{crear} y \emph{evaluar} respectivamente,
todos iguales o inferiores al del objetivo general.

\subsection{Criterios SMART}
\label{subsec:smart}

\textbf{OE1.} \emph{Específico}: compara tres detectores concretos, STA/LTA,
PhaseNet y EQTransformer, sobre un conjunto de datos delimitado.
\emph{Medible}: \emph{recall} y precisión globales y por clase con intervalos
de confianza del 95\,\% obtenidos por \emph{bootstrap}; distribución de IoU;
error de onset; frontera de Pareto de los barridos de parámetros.
\emph{Alcanzable}: ejecutado en su totalidad con inferencia en CPU sobre
modelos preentrenados, sin requerir entrenamiento.
\emph{Relevante}: aborda los aspectos P1 y P2 del problema
(Cuadro~\ref{tab:trazabilidad}). \emph{Temporizado}: cerrado y documentado
para la Etapa 2, el 2 de octubre de 2026.

\textbf{OE2.} \emph{Específico}: define un procedimiento de recorte cuya
extensión depende de la duración detectada, y su transformación a una entrada
de dimensiones compatibles con una red convolucional. \emph{Medible}: IoU del
recorte adaptativo contra el recorte de referencia, comparado con la línea
base de ventana fija y desglosado por clase; fracción de eventos truncados;
cobertura del rango de duraciones de 3,4\,s a 240,8\,s; número de solapamientos
resueltos correctamente. \emph{Alcanzable}: opera sobre las salidas ya
generadas por OE1, sin dependencia de hardware especializado.
\emph{Relevante}: aborda P3 y P1. \emph{Temporizado}: diseñado para la
Etapa 2 e implementado y medido para la Etapa 3, el 30 de octubre de 2026.

\textbf{OE3.} \emph{Específico}: contrasta un conjunto acotado de tres
arquitecturas convolucionales sobre los recortes producidos por OE2.
\emph{Medible}: F1 macro y \emph{recall} por clase con intervalos de confianza
del 95\,\%; matriz de confusión; diferencia de desempeño del mismo modelo
evaluado sobre recorte de referencia y sobre recorte automático; costo de
inferencia en CPU. \emph{Alcanzable}: el entrenamiento se realiza en Google
Colab y la inferencia en CPU local; el conjunto de entrenamiento de 9\,102
eventos etiquetados ya está disponible. \emph{Relevante}: aborda P4 y P5.
\emph{Temporizado}: resultados preliminares en la Etapa 3 y validación
completa en la Etapa 4, el 4 de diciembre de 2026.

\subsection{Matriz de trazabilidad}
\label{subsec:trazabilidad}

El Cuadro~\ref{tab:trazabilidad} vincula cada aspecto del problema
identificado en la Sección~\ref{sec:problema} con el objetivo que lo aborda y
con la medición que permitirá verificar su cumplimiento.

{\footnotesize
\begin{longtable}{@{}p{0.055\textwidth}p{0.245\textwidth}p{0.245\textwidth}p{0.05\textwidth}p{0.245\textwidth}@{}}
\caption{Trazabilidad entre los aspectos del problema, la evidencia que los
respalda, el objetivo que los aborda y su criterio de verificación.}\label{tab:trazabilidad}\\
\toprule
\textbf{Cód.} & \textbf{Aspecto del problema} & \textbf{Evidencia} &
\textbf{OE} & \textbf{Verificación} \\
\midrule
\endfirsthead
\multicolumn{5}{l}{\itshape Cuadro \thetable{} (continuación)}\\
\toprule
\textbf{Cód.} & \textbf{Aspecto del problema} & \textbf{Evidencia} &
\textbf{OE} & \textbf{Verificación} \\
\midrule
\endhead
\midrule
\multicolumn{5}{r}{\itshape Continúa en la página siguiente}\\
\endfoot
\bottomrule
\endlastfoot
P1 & El umbral de detección decide sobre energía y no distingue eventos
volcánicos de no volcánicos &
36\,\% del catálogo continuo corresponde a AV o IC, y los detectores los
encuentran con desempeño significativamente distinto por clase &
OE1, OE2 &
Separabilidad del estadístico de detección entre clases de interés y de no
interés \\
\addlinespace
P2 & La elección del detector carece de respaldo empírico local &
Caídas de \emph{recall} de 13--56\,\% fuera del dominio de entrenamiento
\parencite{jiang2021comparacion} &
OE1 &
\emph{Recall}, precisión e IoU por detector y por clase con IC del 95\,\% \\
\addlinespace
P3 & La ventana de análisis de longitud fija trunca eventos &
Duración mediana de TR de 87,8\,s frente a una ventana de 81,92\,s; rango de
3,4 a 240,8\,s &
OE2 &
Fracción de eventos truncados e IoU del recorte adaptativo contra la línea
base de ventana fija \\
\addlinespace
P4 & No se ha medido la pérdida de desempeño al pasar de recortes de
referencia a recortes automáticos &
La totalidad de los trabajos revisados entrena y evalúa sobre segmentaciones
de referencia &
OE3 &
Diferencia de F1 macro del mismo modelo sobre ambos tipos de recorte \\
\addlinespace
P5 & El clasificador vigente no cubre las clases no volcánicas y las resuelve
por rechazo estadístico &
El módulo OOD rechaza entre 59\,\% y 100\,\% de lo detectado, existiendo 1\,782
ejemplos etiquetados de AV e IC &
OE3 &
\emph{Recall} de la clase de rechazo y tasa de entradas fuera de dominio
atribuidas a la etapa de detección \\
\bottomrule
\end{longtable}
}

\section{Alcances y limitaciones}
\label{sec:alcances}

\begin{itemize}
    \item Esta tesis \textbf{no} diseña arquitecturas de red nuevas. La
    contribución consiste en configurar, adaptar y evaluar componentes
    existentes sobre datos sismo-volcánicos del NVChVC.

    \item Los detectores basados en aprendizaje profundo se emplean
    \textbf{preentrenados}, sin reentrenamiento ni ajuste fino sobre datos
    volcánicos chilenos. Esta decisión se declara como limitación y se propone
    como trabajo futuro.

    \item La detección opera a nivel de \textbf{canal individual}. La
    asociación por coincidencia entre estaciones queda fuera del alcance y se
    señala como la vía natural para reducir falsos positivos en trabajo
    posterior.

    \item El trabajo se restringe al \textbf{Complejo Volcánico Nevados de
    Chillán}. No se evalúa la generalización a otros volcanes ni a otras redes
    de estaciones.

    \item La evaluación se realiza en modalidad \textbf{fuera de línea}, sobre
    registro histórico. El despliegue en producción, la ingesta en tiempo real
    y la integración con el flujo operativo del observatorio quedan excluidos.

    \item Toda inferencia se ejecuta en \textbf{CPU}, dado que el entorno de
    desarrollo no dispone de aceleración por hardware. El entrenamiento se
    realiza en un entorno externo con GPU, manteniendo compatibilidad de
    inferencia en CPU mediante el mecanismo estándar de PyTorch.

    \item El catálogo de referencia se toma como \textbf{verdad de terreno} sin
    cuestionar su incertidumbre. No se realiza un estudio de concordancia entre
    analistas.
\end{itemize}

%% ======================================================================
%% Se inserta al final de cap1.tex con:  \input{cap1_requerimientos}
%% No lleva \chapter: sus secciones pertenecen al Capitulo 1.
%% El detalle tabular vive en anexo_requerimientos.tex
%% ======================================================================

\section{Análisis de requerimientos}
\label{sec:requerimientos}

Los requerimientos de este trabajo no se postulan a priori: se derivan de tres
fuentes diferenciables. La primera son las mediciones propias obtenidas en la
evaluación de detectores, ya cerrada. La segunda son las restricciones del
entorno de trabajo y de los datos disponibles. La tercera son las decisiones
metodológicas acordadas con el profesor guía. Cada requerimiento declara su
origen y el criterio con el que se comprobará su cumplimiento, de modo que
ninguno quede formulado como una aspiración no verificable. Cuando el origen es
una medición, se indica el valor observado; cuando es un supuesto todavía no
verificado, se declara como tal y no se lo trata como evidencia.

\subsection{Alcance del sistema}
\label{sec:alcance}

El sistema recibe registro sísmico continuo y entrega, para cada evento
detectado, una etiqueta de clase con su puntaje asociado. Se organiza en cuatro
bloques encadenados. El bloque A ingiere la traza multicanal a 100~Hz y la
acondiciona, aplicando la estrategia de relleno de componentes que corresponda al
detector. El bloque B detecta eventos y entrega intervalos candidatos con
instante de inicio, término y puntaje. El bloque C recorta cada candidato con
duración variable y produce una representación espectral de dimensiones fijas. El
bloque D clasifica esa representación en cuatro salidas: VT, LP, TR y rechazo.

El bloque B ya fue estudiado empíricamente y su evaluación está cerrada; sus
resultados son insumo de los bloques C y D, no materia pendiente. Los bloques C
y D constituyen el desarrollo propio de este trabajo.

La identificación de las partes interesadas y la evidencia documental que
sustenta la relevancia del problema se presentaron en la
Sección~\ref{subsec:interesados}.

\subsection{Evidencia empírica que origina los requerimientos}
\label{sec:evidencia}

La evaluación de detectores entregó nueve hallazgos que condicionan el diseño de
los bloques posteriores, resumidos en el Cuadro~\ref{tab:evidencia}. Todos
corresponden a mediciones sobre la traza continua de 10,18~horas con su catálogo
de 203 eventos.

{\small
\begin{longtable}{>{\raggedright\arraybackslash}p{0.9cm}>{\raggedright\arraybackslash}p{5.3cm}>{\raggedright\arraybackslash}p{5.3cm}}
\caption{Hallazgos medidos en la evaluación de detectores y su implicación de diseño.}\label{tab:evidencia}\\
\toprule
\textbf{ID} & \textbf{Hallazgo medido} & \textbf{Implicación de diseño} \\
\midrule
\endfirsthead
\multicolumn{3}{l}{\itshape Cuadro \thetable{} (continuación)}\\
\toprule
\textbf{ID} & \textbf{Hallazgo medido} & \textbf{Implicación de diseño} \\
\midrule
\endhead
\midrule
\multicolumn{3}{r}{\itshape Continúa en la página siguiente}\\
\endfoot
\bottomrule
\endlastfoot
E1 & El recall global de los tres detectores es estadísticamente indistinguible:
0,542; 0,527 y 0,581, con intervalos de confianza del 95\,\% solapados &
El detector no puede seleccionarse por su recall global \\
\addlinespace
E2 & Por clase existen diferencias significativas de signo opuesto que se cancelan
al promediar: un detector aventaja en LP y TR, el otro en VT, AV e IC &
Toda métrica debe reportarse desagregada por clase \\
\addlinespace
E3 & En calidad de delimitación, STA/LTA supera a EQTransformer de forma robusta
frente al umbral de solapamiento: IoU de 0,563 contra 0,479 &
La variabilidad del recorte entre detectores es un hecho medido; el clasificador
debe absorberla \\
\addlinespace
E4 & PhaseNet no delimita eventos: entrega arribos, no intervalos &
Se requiere una regla explícita de expansión de arribo a intervalo \\
\addlinespace
E5 & La mediana de duración de TR es de 87,8~s y excede la ventana fija de
81,92~s; el rango observado va de 3,4~s a 240,8~s &
El recorte debe ser adaptativo: la ventana fija trunca la clase más extensa \\
\addlinespace
E6 & 73 de los 203 eventos del catálogo continuo, un 36\,\%, son de clases no
volcánicas, y los detectores las encuentran con buen desempeño &
El rechazo no puede ocurrir en la detección; debe resolverse en la clasificación \\
\addlinespace
E7 & De 106 falsas alarmas analizadas, un 20,8\,\% son fragmentos de eventos
catalogados, un 35,8\,\% son adyacentes y un 43,4\,\% aisladas &
Las falsas alarmas son material de entrenamiento aprovechable como negativos
duros \\
\addlinespace
E8 & La estrategia de relleno de componentes altera el techo de recall y su efecto
depende del modelo: uno pasa de 0,527 a 0,946, el otro se comporta de forma
inversa &
El preprocesamiento es un parámetro del sistema, no una constante \\
\addlinespace
E9 & El error de onset se sitúa entre 1,7~s y 2,0~s en los modelos profundos y en
2,4~s en STA/LTA &
El recorte debe incorporar margen de contexto previo \\
\bottomrule
\end{longtable}
}

Conviene explicitar la tensión de diseño que surge de E5. Una red convolucional
exige entrada de dimensiones fijas, mientras que las duraciones observadas
abarcan casi dos órdenes de magnitud. Reconciliar ambas exigencias sin destruir
la información de duración, que es discriminante entre clases, es el problema
central del bloque C y la razón de ser del segundo objetivo específico.

\subsection{Evaluación de alternativas}
\label{sec:alternativas}

Cada decisión de diseño se tomó contrastando alternativas viables. El
Cuadro~\ref{tab:alternativas} las resume; el desarrollo del criterio decisivo
en cada caso está en el Anexo~\ref{sec:anexo-alternativas}.

{\small
\begin{longtable}{>{\raggedright\arraybackslash}p{0.8cm}>{\raggedright\arraybackslash}p{3.0cm}>{\raggedright\arraybackslash}p{3.6cm}>{\raggedright\arraybackslash}p{3.7cm}}
\caption{Alternativas consideradas y criterio de selección.}\label{tab:alternativas}\\
\toprule
\textbf{ID} & \textbf{Decisión} & \textbf{Alternativas} & \textbf{Selección y criterio} \\
\midrule
\endfirsthead
\multicolumn{4}{l}{\itshape Cuadro \thetable{} (continuación)}\\
\toprule
\textbf{ID} & \textbf{Decisión} & \textbf{Alternativas} & \textbf{Selección y criterio} \\
\midrule
\endhead
\midrule
\multicolumn{4}{r}{\itshape Continúa en la página siguiente}\\
\endfoot
\bottomrule
\endlastfoot
D1 & Elección del detector &
Detector único elegido por desempeño; conservar varios &
Conservar varios: el recorte variable es el objeto de estudio, no un defecto a
eliminar \\
\addlinespace
D2 & Manejo de clases no volcánicas &
Umbral en detección; detección OOD posterior; clase de rechazo supervisada &
Rechazo supervisado: hay 1.782 ejemplos etiquetados y los detectores sí
encuentran esas clases \\
\addlinespace
D3 & Estrategia de recorte &
Ventana fija de 81,92~s; ventana fija mayor; recorte adaptativo &
Recorte adaptativo: la ventana fija trunca TR y una mayor degrada la relación
entre evento y contexto \\
\addlinespace
D4 & Representación de entrada &
Forma de onda cruda; espectrograma; escalograma &
Representación espectral: mantiene comparabilidad con la línea base publicada \\
\addlinespace
D5 & Arquitecturas a comparar &
Una sola arquitectura; comparación de varias &
Tres arquitecturas: la elección debe quedar respaldada y no ser arbitraria \\
\addlinespace
D6 & Uso de detectores profundos &
Preentrenados; ajuste fino con datos locales &
Preentrenados: el ajuste fino consumiría el único conjunto con delimitación de
referencia \\
\addlinespace
D7 & Nivel de operación &
Canal individual; coincidencia multiestación &
Canal individual: la asociación robusta excede el presupuesto de esta etapa \\
\bottomrule
\end{longtable}
}

\subsection{Síntesis de requerimientos}
\label{sec:sintesis-req}

El sistema queda especificado por veinte requerimientos funcionales y ocho no
funcionales. El Cuadro~\ref{tab:sintesis} resume su distribución; el detalle de
cada uno, con su origen y su criterio de verificación, está en el
Anexo~\ref{sec:anexo-rf}.

{\small
\begin{longtable}{>{\raggedright\arraybackslash}p{3.4cm}>{\raggedright\arraybackslash}p{2.0cm}>{\raggedright\arraybackslash}p{5.8cm}}
\caption{Distribución de los requerimientos por bloque funcional.}\label{tab:sintesis}\\
\toprule
\textbf{Bloque} & \textbf{Requerim.} & \textbf{Exigencia central} \\
\midrule
\endfirsthead
\multicolumn{3}{l}{\itshape Cuadro \thetable{} (continuación)}\\
\toprule
\textbf{Bloque} & \textbf{Requerim.} & \textbf{Exigencia central} \\
\midrule
\endhead
\midrule
\multicolumn{3}{r}{\itshape Continúa en la página siguiente}\\
\endfoot
\bottomrule
\endlastfoot
A. Ingesta y preprocesamiento & RF-01 a RF-05 &
Lectura correcta, canales utilizables, relleno declarado por modelo y partición
libre de superposición \\
\addlinespace
B. Detección & RF-06 a RF-08 &
Interfaz común de detector, expansión de arribo a intervalo y métricas por clase \\
\addlinespace
C. Recorte y representación & RF-09 a RF-13 &
Duración adaptativa, política para eventos extensos, margen de contexto,
solapamiento y entrada de tamaño fijo \\
\addlinespace
D. Clasificación & RF-14 a RF-17 &
Cuatro salidas, entrenamiento del rechazo, desbalance y comparación de
arquitecturas \\
\addlinespace
Evaluación integrada & RF-18 a RF-20 &
Comparación pareada, intervalos de confianza y persistencia reproducible \\
\addlinespace
Calidad del producto & RNF-01 a RNF-08 &
Inferencia en CPU, entrenamiento reanudable, reproducibilidad, modularidad,
eficiencia y salidas revisables \\
\bottomrule
\end{longtable}
}

\subsection{Restricciones y supuestos}
\label{sec:restricciones}

Las restricciones acotan el espacio de soluciones admisibles y no son negociables
dentro de este alcance. El equipo de trabajo carece de aceleración por GPU, lo que
obliga a entrenar en un entorno remoto y a diseñar la inferencia para CPU. El
conjunto de datos es fijo y no se solicitarán registros adicionales. Solo un
subconjunto de las estaciones presenta señal utilizable en la traza continua. Los
nombres de archivo del conjunto de entrenamiento no incluyen el año, lo que impide
descartar por metadatos la superposición con la traza de evaluación. Los
detectores profundos se emplean con pesos preentrenados. El calendario académico
fija cuatro hitos de evaluación y condiciona la profundidad alcanzable en cada
frente.

La traza continua de evaluación no es un registro ininterrumpido. Al inspeccionar
la fila de marca de tiempo del archivo se identifican dos discontinuidades, en las
muestras 2.160.000 y 3.240.001, que separan tres segmentos procedentes de fechas
distintas y no consecutivas: seis horas del 9 de febrero de 2018, tres horas del
14 de enero de 2020 y una hora y once minutos del 10 de enero de 2017, todas en
UTC. Esta estructura no aparece descrita en la documentación del conjunto y se
determinó de forma directa sobre el archivo. Tiene tres consecuencias: obliga a
excluir las guardas de unión entre segmentos al construir el conjunto de
evaluación, sostiene la partición por bloque temporal como mitigación del riesgo
RG-01, e impide descartar por metadatos la superposición con el conjunto de
entrenamiento, cuyos nombres registran día, mes y hora pero no año.

\subsection{Riesgos y mitigaciones}
\label{sec:riesgos}

Los riesgos que pueden comprometer la validez de los resultados o el
cumplimiento del calendario se detallan en el Anexo~\ref{sec:anexo-riesgos}.
Dos son de severidad alta. El primero es la posible superposición entre el
conjunto de entrenamiento y la traza de evaluación: hay 31 archivos cuya fecha
parcial cae dentro de los bloques de la traza, pero los nombres carecen de año
y los conteos son compatibles con el azar, de modo que no está probada ni
descartada. Se mitiga con partición por bloque temporal y verificación directa
por correlación de forma de onda antes de reportar resultados. El segundo es el
truncamiento sistemático de los eventos más extensos por efecto de la ventana
fija, que se mitiga con el recorte adaptativo y su política explícita para
eventos que exceden la ventana nominal. La severidad corresponde a un juicio
propio y se declara como tal, no como una estimación cuantitativa.

\subsection{Criterios de aceptación}
\label{sec:aceptacion}

Los criterios se formulan como comparaciones contra una referencia medible y no
como umbrales absolutos, porque fijar un umbral numérico sin medición previa que
lo respalde sería arbitrario.

\begin{itemize}
\item \textbf{CA-01.} El clasificador entrenado sobre recortes adaptativos debe
superar, con intervalos de confianza que no se solapen, a un clasificador de
referencia entrenado sobre la ventana fija actual, en las métricas por clase.
\item \textbf{CA-02.} La degradación entre recorte de referencia y recorte
automático debe quedar cuantificada con su intervalo de confianza. No se fija un
máximo admisible: la magnitud de esa degradación es el resultado que el tercer
objetivo específico busca establecer.
\item \textbf{CA-03.} La clase de rechazo debe reducir la proporción de
detecciones no volcánicas propagadas como eventos volcánicos, respecto de la
configuración sin rechazo.
\item \textbf{CA-04.} Toda comparación entre arquitecturas debe realizarse bajo
partición, semilla y presupuesto idénticos, y reportarse por clase.
\item \textbf{CA-05.} La cadena completa debe ejecutarse sobre la traza de
evaluación en CPU y reportar su tiempo de procesamiento.
\end{itemize}

\subsection{Matriz de trazabilidad}
\label{sec:trazabilidad-req}

El Cuadro~\ref{tab:trazabilidad} enlaza cada aspecto del problema con el objetivo
específico que lo aborda y con los requerimientos que lo materializan. Verifica
que ningún aspecto quede sin requerimiento asociado y permite comprobar en la
etapa final que cada objetivo fue cubierto.

{\small
\begin{longtable}{>{\raggedright\arraybackslash}p{0.8cm}>{\raggedright\arraybackslash}p{4.6cm}>{\raggedright\arraybackslash}p{1.3cm}>{\raggedright\arraybackslash}p{4.8cm}}
\caption{Trazabilidad de requerimientos.}\label{tab:trazabilidad-req}\\
\toprule
\textbf{ID} & \textbf{Aspecto del problema} & \textbf{Objetivo} & \textbf{Requerimientos} \\
\midrule
\endfirsthead
\multicolumn{4}{l}{\itshape Cuadro \thetable{} (continuación)}\\
\toprule
\textbf{ID} & \textbf{Aspecto del problema} & \textbf{Objetivo} & \textbf{Requerimientos} \\
\midrule
\endhead
\midrule
\multicolumn{4}{r}{\itshape Continúa en la página siguiente}\\
\endfoot
\bottomrule
\endlastfoot
P1 & El umbral decide sobre energía y no sobre clase & OE1, OE2 &
RF-06, RF-07, RF-08, RF-13, RF-14 \\
\addlinespace
P2 & La elección del detector carece de respaldo empírico local & OE1 &
RF-03, RF-06, RF-08, RF-19 \\
\addlinespace
P3 & La ventana fija trunca los eventos más extensos & OE2 &
RF-09, RF-10, RF-11, RF-12, RF-13 \\
\addlinespace
P4 & No se ha medido la degradación entre recorte de referencia y automático &
OE3 & RF-18, RF-19, RF-20 \\
\addlinespace
P5 & El clasificador vigente no cubre las clases no volcánicas & OE3 &
RF-14, RF-15, RF-16, RF-17 \\
\bottomrule
\end{longtable}
}

Los requerimientos RF-01, RF-02, RF-04 y RF-05 no se asocian a un aspecto
específico porque son condiciones habilitantes: sin lectura correcta de los datos
y sin una partición libre de superposición, ninguna medición posterior sería
interpretable.

\section{Planificación del proyecto}
\label{sec:planificacion}

\subsection{Capacidad y método de estimación}
\label{sec:capacidad}

La capacidad de trabajo se estima en 20~horas semanales, valor medio de una
disponibilidad declarada entre 15 y 25~horas, sobre la que se aplica un factor de
disponibilidad efectiva de 0,8 para absorber interrupciones y tareas
administrativas. Resultan 16~horas comprometibles por semana. El período hasta la
entrega final comprende 13,1 semanas repartidas en tres sprints de cuatro semanas,
con una capacidad total cercana a 208~horas.

Las estimaciones por tarea provienen de juicio propio anclado en trabajo análogo
ya ejecutado durante la evaluación de detectores. Corresponde declarar la
limitación: no se llevó registro del esfuerzo real invertido en aquellas tareas,
por lo que el anclaje es cualitativo y las cifras son estimaciones, no
mediciones. Como mitigación, registraré el esfuerzo real de cada tarea a medida
que se complete, de modo que los sprints siguientes puedan recalibrarse con datos
propios.

\subsection{Épicas}
\label{sec:epicas}

El backlog se organiza en seis épicas. \textbf{EP-1, validez experimental},
garantiza que ningún resultado esté contaminado por superposición entre conjuntos
ni por supuestos no verificados. \textbf{EP-2, cadena de detección operativa},
permite conectar cualquier detector mediante una interfaz única. \textbf{EP-3,
recorte adaptativo y representación}, resuelve el paso de duración variable a
entrada fija y responde a OE2. \textbf{EP-4, clasificador de cuatro salidas},
distingue las clases volcánicas y rechaza el resto. \textbf{EP-5, evaluación
integrada}, mide la degradación entre recorte de referencia y automático, y
responde a OE3. \textbf{EP-6, documentación y cierre}, asegura que cada resultado
quede escrito, versionado y reproducible.

Cada épica se descompone en historias de usuario y estas en tareas técnicas
estimadas y priorizadas. El desglose completo, con dependencias y criterios de
aceptación, está en el Anexo~\ref{sec:anexo-backlog}.

\subsection{Distribución por sprints}
\label{sec:sprints}

{\small
\begin{longtable}{>{\raggedright\arraybackslash}p{1.5cm}>{\raggedright\arraybackslash}p{1.8cm}>{\raggedright\arraybackslash}p{5.2cm}>{\raggedright\arraybackslash}p{1.3cm}>{\raggedright\arraybackslash}p{1.3cm}}
\caption{Distribución del backlog en sprints de cuatro semanas.}\label{tab:sprints}\\
\toprule
\textbf{Sprint} & \textbf{Cierre} & \textbf{Foco} & \textbf{Compr.} & \textbf{Holgura} \\
\midrule
\endfirsthead
\multicolumn{5}{l}{\itshape Cuadro \thetable{} (continuación)}\\
\toprule
\textbf{Sprint} & \textbf{Cierre} & \textbf{Foco} & \textbf{Compr.} & \textbf{Holgura} \\
\midrule
\endhead
\midrule
\multicolumn{5}{r}{\itshape Continúa en la página siguiente}\\
\endfoot
\bottomrule
\endlastfoot
S1 & 2 de octubre & Validez experimental y diseño de la cadena: verificación de
superposición, diagramas de arquitectura y especificación de interfaces & 58 h & 6 h \\
\addlinespace
S2 & 30 de octubre & Implementación del recorte adaptativo, de la representación
espectral y del conjunto de entrenamiento de cuatro clases & 58 h & 6 h \\
\addlinespace
S3 & 4 de diciembre & Comparación de arquitecturas, evaluación integrada y
redacción del informe final & 72 h & 8 h \\
\bottomrule
\end{longtable}
}

La distribución responde a las dependencias técnicas entre tareas y a la
naturaleza de cada etapa de evaluación: el primer sprint concentra el diseño, el
segundo la implementación y el tercero la validación y el cierre documental. Ante
una desviación, aplicaré la siguiente política: si una tarea supera su estimación
en más del 50\,\%, se detiene su ejecución y se reevalúa el alcance en lugar de
absorber la holgura de forma silenciosa, de modo que la desviación se convierta en
información para replanificar y no en una deuda invisible.